\begin{center}
        \begin{tikzpicture}[
        scale=1,
        point/.style={circle, fill, inner sep=1.5pt}
    ]

    % Punkte
    \coordinate (C) at (0,0);
    \coordinate (B) at (3,0);
    \coordinate (A) at (1.5,2.25);
    \coordinate (D) at (1.5,0);
    \coordinate (E) at (2.25,1.125);

    % Dreieck
    \draw[thick]
        (B) -- node[below] {$a$}
        (C) -- node[left, sloped, above] {$b$}
        (A) -- node[right, sloped, above] {$c$}
        cycle;

    % Lotfuß
    \draw[thick, red!40] (A) -- node[midway, right, text=red!40] {$h'$} (D);
    \draw[thick, green!40] (C) -- node[midway, above, text=green] {$h$} (E);

    % Punkte markieren
    \node[point, label=below left:$C$] at (C) {};
    \node[point, label=below right:$B$] at (B) {};
    \node[point, label=above:$A$] at (A) {};
    \fill[red!40] (D) circle (1.5pt) node[above right, text=red!40] {$D$};
    \fill[green] (E) circle (1.5pt) node[below, text=green] {$E$};

    % Winkel gamma bei C
    \pic[
        draw,
        angle radius=3mm,
        angle eccentricity=1.4,
        pic text={$\gamma$}
    ] {angle=B--C--A};

    % Winkel beta bei B
    \pic[
        draw,
        angle radius=3mm,
        angle eccentricity=1.4,
        pic text={$\beta$}
    ] {angle=A--B--C};

    % Winkel alpha bei A
    \pic[
        draw,
        angle radius=3mm,
        angle eccentricity=1.4,
        pic text={$\alpha$}
    ] {angle=C--A--B};

    \pic[
        draw,
        angle radius=2mm,
        angle eccentricity=0.5,
        pic text={\fontsize{3}{3}\selectfont$\bullet$}
    ] {angle=A--D--C};

    \pic[
        draw,
        angle radius=2mm,
        angle eccentricity=0.5,
        pic text={\fontsize{3}{3}\selectfont$\bullet$}
    ] {angle=A--E--C};
            
    \end{tikzpicture}